\section{Лекция 21}
\subsection{Обобщение поверхностных интегралов на кусочно-гладкую поверхность}
\begin{definition}
    Кусочно-гладкой поверхностью $\Sigma$ в $\R^3$ назовем конечный набор гладких поверхностей $\{\Sigma_j\}$, где $\Sigma_j$ с разными индексами либо не пересекаются, либо пересекаются по гладкой дуге границы. Любые три разных $\Sigma_j$ могут пересекаться по одной точке. Из любой $\Sigma_j$ можно перейти в любую $\Sigma_j$ по цепочке соприкасающихся поверхностей. 
    Интеграл первого рода определяется так:
    \[\iint\limits_{\Sigma}f\ dS =\sum\limits_{j=1}^{n}\iint\limits_{\Sigma_j}f\ dS\]
    Интеграл второго рода определяется также: (по аддитивности)
\end{definition}
\begin{definition}
    Ориентированной кусочно-гладкой поверхностью называется кусочно гладкая поверхность у которой линии склейки проходятся в разном направлении.
\end{definition}
\begin{comm}
    Таким образом на кусочно-гладкую поверхность можно перенести формулу Стокса.
\end{comm}
\begin{theorem}[Формула Гаусса-Остроградского] $\\$
    Пусть $G$ --- выпуклая область в $\R^3$ с кусочно-гладкой границей $\partial G$, которая ориентированна полем внешних нормалей. На $\overline{G}$ определено $\mathcal{C}^1$ векторное поле $(P,Q,R)$. 
    \[\iint\limits_{\partial G} Pdydz+ Qdzdx+ Rdxdy=\iiint\limits_{\overline{G}}(P'_x+Q'_y+R'_z)\ dxdydz\]
\end{theorem}
\begin{proof}
    Докажем для одного слагаемого.\\
    Спроецируем $G$ на плоскость $XOY$. $(x,y)\in \overline{\Omega},\ \psi_1(x,y)\leq z\leq \psi_2(x,y)$
    \[\iiint\limits_{\overline{G}}R'_z\ dxdydz=\iint\limits_{\overline{\Omega}}\left(\ \int\limits_{\psi_1(x,y)}^{\psi_2(x,y)}R'_z(x,y,z)\ dz\right)dxdy\]
    \[\int\limits_{\psi_1(x,y)}^{\psi_2(x,y)}R'_z(x,y,z)\ dz=R(x,y,\psi_2(x,y))-R(x,y,\psi_1)\]
    там как то поделили на сигма 1 2 и 3
    \[\iint\limits_{\overline{\Sigma_2}}R(x,y,z)\ dxdy=\iint\limits{\overline{\Omega}}R(x,y,\psi_2(x,y))\ dxdy\]
    Получается
    \[\iint\limits_{\Sigma_2^+} R\ dxdy+\iint\limits_{\Sigma_1^-}R\ dxdy+\iint\limits_{\Sigma_3}R\ dxdy\]
\end{proof}
\subsection{Теоремы Гульдина}
Зададим плотность $\rho(x,y)$ в каждой точке $\overline{\Omega}$.
\[x_0=\frac{\iint\limits_{\overline{\Omega}}x\cdot \rho(x,y)\ dxdy}{\iint\limits_{\overline{\Omega}}\rho(x,y)\ dxdy}\]
\begin{theorem}[Первая теорема Гульдина] $\\$ % ЗАМЕЧАНИЕ НА НЕКСТ ЛЕКЦИИ ДАЛ: ВМЕСТО НАТУРАЛЬНОЙ ПАРАМЕТРИЗАЦИИ МОЖНО ИСПОЛЬЗОВАТЬ ЛЮБУЮ (можно и натуральную то тогда это работается только для кусочно гладких кривых)
    Пусть дано направление $l$ и $\gamma$ --- кусочно-гладкая простая замкнутая кривая. Определим $\overline{G}$ --- тело вращения $\overline{\Omega}$ вокруг $l$. Тогда $\Sigma=\partial G$ --- поверхность вращения. Площадь $\Sigma$ равна длине $\gamma$ умноженой на длину пути центра масс кривой $\gamma$ при вращении вокруг $l$.
\end{theorem}
\begin{proof}
    Параметризуем поверхность
    \[\begin{cases}
        x=\widehat{x}(s)\cdot \cos{\phi},\\
        y=\widehat{x}(s)\cdot \sin{\phi},\\
        z=\widehat{z}(s).
    \end{cases}\]
    \[r'_s=(\widehat{x}'(s)\cdot \cos{\phi}, \widehat{x}'(s)\cdot \sin{\phi}, \widehat{z}'(s))\]
    \[r'_{\phi}=(-\widehat{x}(s)\cdot \sin{\phi},\widehat{x}\cdot \cos{\phi}, 0)\]
    Площадь сигма это
    \[\int\limits_{0}^{2\pi}\int\limits_{0}^{|\gamma|}\widehat{x}(s)\ dsd\phi=2\pi\cdot \int\limits_{0}^{|\gamma|}\]
    \[\sqrt{EG-F^2}=\widehat{x}(s)\]
    но 
    \[x_0=\frac{\int\limits_{\gamma}x\ ds}{\int\limits_{\gamma}1\ ds}\]
\end{proof}
\begin{theorem}[Вторая теорема Гульдина] $\\$
    Объем $G$ равен площади $\Omega$ умноженую на длину пути центра масс области $\overline{\Omega}$ при вращении вокруг $l$.
\end{theorem}
\begin{proof}
    Параметризуем. Пусть $(u,v)\in \overline{\Omega}$.
    \[\begin{cases}
        x=u\cdot \cos{\phi},\\
        y=u\cdot \sin{\phi},\\
        z=v.
    \end{cases}\]text
    \[\int\limits_{0}^{2\pi}\iint\limits_{\overline{\Omega}}u\ dudvd\phi=2\pi\cdot \iint\limits_{\overline{\Omega}}u\ dudv\]
    при этом
    \[u_0=\cfrac{\iint\limits_{\overline{\Omega}}u\ dudv}{\iint\limits_{\overline{\Omega}}1\ dudv}\]
\end{proof}
%пример с тором
\newpage